% docs/.paper/sec_appendix_primitives.tex
\section{Methodology Primitives}
\label{app:primitives}

This appendix details the foundational primitives derived for repository-scale software migration. Marker legend: \texttt{[MEASURED]} --- directly observed in primary literature; \texttt{[AUTHOR\_CLAIM]} --- authors stated, no independent replication; \texttt{[INFERRED]} --- derived reasoning across multiple sources; \texttt{[HYPOTHESIS]} --- proposed but unvalidated.

\begin{enumerate}

\item \textbf{Reverse Topological Translation Ordering.}
Bottom-up implementation plan specifying target module order (ReCodeAgent [P01: Sec 3.3, p. 5]) / reverse-topological method translation (AlphaTrans [P02 \S6 \P1]).

\item \textbf{Back-Edge-Conditioned Reverse-Topological Scheduling.}
Build a directed call graph, detect cycles, remove back edges to obtain a DAG, compute a reverse-topological linearization, and translate leaves first (AlphaTrans [P02 \S6 \P1]).

\item \textbf{Target Skeleton-First Generation.}
Before method bodies, emit target-language interface stubs, type declarations, imports, and signature-only functions for every source module (ReCodeAgent [P01 Sec 3.3.3]; AlphaTrans [P02 \S5.2]).

\item \textbf{SpecMiner-Style Dynamic Invariant Recovery.}
Instrument source binaries at runtime to record allocation sizes, nullability of pointer targets, aliasing relationships, lifetime ranges, and branch coverage under representative inputs (Syzygy [P58 \S3.2, p. 4]).

\item \textbf{Test Suite Co-Translation \& Synthesis.}
Translate source-language unit tests into the target language as part of the translation pipeline; when source tests are insufficient, generate additional tests for uncovered functions (ReCodeAgent [P01 Sec 3.5.2; Alg 1 lines 14--21]; AlphaTrans [P02 \S6 --- testCheck]; AdvTestGen [P25]).

\item \textbf{Multi-Stage Build/Test Feedback Repair.}
After initial translation, run a 3-tier validation cascade (AST/syntax check, target compiler check, translated test suite execution) and feed diagnostics back to the translator up to a fixed budget (ReCodeAgent [P01 Alg 1]; AlphaTrans [P02 \S6 Alg 2]).

\item \textbf{Multi-Agent Verdict Validation.}
After translating, deploy multiple LLM verdict agents to independently judge whether the target translation is semantically equivalent to the source; on disagreement, prompt the translator to repair (MatchFixAgent [P08 \S3]).

\item \textbf{State-Grounded Mock-Based In-Isolation Validation.}
For each module under translation, generate a mock of its dependencies using state-equivalent stubs and validate the module against the mock in isolation, then assemble and validate the full system (TRAM [P10 \S3]).

\item \textbf{Tri-Representation Hybrid Code Graph.}
A unified graph combining the Abstract Syntax Tree (AST, Tree-sitter structural parse), Code Property Graph (CPG, Yamaguchi AST $\cup$ CFG $\cup$ PDG), and System Dependence Graph (SDG, Horwitz/Reps/Binkley inter-procedural PDG extension). RepoGraph [P62] grounds the typed code-graph layer empirically with a $+32.8\%$ average relative resolve-rate lift across four SWE-agent baselines and a $+93\%$ Code Match EM lift on CrossCodeEval with a small open model.

\item \textbf{Feature-Mapping \& Type-Compatibility Validation.}
Maintain an explicit mapping table from source-language idioms to target-language equivalents (e.g., Go interfaces $\leftrightarrow$ Rust traits; Java checked exceptions $\leftrightarrow$ Rust \texttt{Result}) and validate every target symbol against its source counterpart (Oxidizer [P05 \S3]; RustRepoTrans [P07 \S3]).

\item \textbf{Implementation-Agnostic Testing.}
Validate target behavior using black-box I/O test vectors rather than asserting on the structure of the translated code (RepoMod-Bench [P11 \S3]).

\item \textbf{Wasm-Based Reference Execution Oracle.}
Compile the original source code to WebAssembly; run the same input through the Wasm-compiled source and the target-language implementation; compare observable outputs (VERT [P12 \S3]).

\item \textbf{Adversarial Test-Weakening Guard.}
During repair iterations, AST-validate that the translator has not weakened tests (deleted assertions, widened tolerances, removed negative cases, relaxed catch blocks) (AdvTestGen [P25: ISSTA 2024]; MatchFixAgent [P08 \S3]).

\item \textbf{Software-Archaeology Stage.}
A pre-translation multi-step comprehension pass: boundary extraction, time-capsule execution harness, DRSpace \& churn bisection, naming \& type forensics, and test \& documentation archaeology (pp-besm dev.to playbook; AgentPatterns.ai Legacy Code Archaeology; Rajlich; M\"uller; Foltz; Baldwin \& Clark; Kazman \& Cai).

\item \textbf{Evidence-First Adaptation Pattern.}
Before adapting external code into a target system, produce four artifacts: (1) target-preserving spec, (2) contract-before-code, (3) provenance-aware diff, (4) untrusted-source model (Reeper [S12 \S2]).

\item \textbf{Spec-Driven Development Lifecycle.}
Strict sequence: constitution $\rightarrow$ specify $\rightarrow$ plan $\rightarrow$ tasks $\rightarrow$ implement $\rightarrow$ converge; each phase produces a reviewable artifact and transitions are gated (spec-kit [S13]).

\item \textbf{Asynchronous Software Engineering Agent Architecture.}
Multiple agents operate concurrently over a shared memory blackboard; agents wake on blackboard events, claim work units, and post results; a central scheduler arbitrates conflicts (CAID [P54 \S3]).

\item \textbf{Jaccard-Coupling Architecture Recovery.}
Compute Jaccard similarity of commit neighborhoods between modules to detect hidden coupling invisible to static import edges (MSR4SA [P52 \S4]).

\item \textbf{Design Rule Hierarchy Partitioning.}
Classify every architectural element by stability layer (L1 interfaces, L2 subsystems, L3 leaves) and apply the load-bearing concentration check (Kazman et al.\ [P51 \S4]).

\item \textbf{Concept Assignment and Redocumentation.}
Before translation, perform concept-locator analysis: identify what concepts the code implements, where they live, and how they relate; output a concept $\rightarrow$ location map (Rajlich [P45 \S2]).

\item \textbf{Migration Strategy Selection.}
For each migration target, classify into one of five strategies: Ignore, COTS, Cold Turkey, Integrate-in-Place, or Chicken Little, plus the LIS (Legacy Integration Strategy) decomposability assessment (M\"uller [P46 slides 18--25]).

\item \textbf{Four Phases of Comprehension.}
Apply the DR.\ JONES model (Decomposition, Recognition, Organization, Navigation, Explanation, Search) with empirically grounded cognitive observations: linear traversal preferred, recency bias, breadth before depth, and hotspot concentration (Foltz [P47 \S3]).

\item \textbf{Chunked Translation.}
Replace single-shot translation with a per-fragment dispatcher that groups planning fragments by topological order, issues one LLM call per source file or per $\sim$500-LoC chunk, and maintains a persistent state summary (emitted modules, public symbols, type aliases, use statements) between calls.

\item \textbf{SOP-Anchored Role-Artifact Schema.}
Each agent emits a strictly-typed intermediate artifact (JSON/markdown with required fields and cardinality constraints); parse-or-retry gates the next agent; each schema carries a version field for forward-compatible migrations (MetaGPT [P42] \S3.2.2).

\item \textbf{Communicative-De-hallucination Role-Flip Gate.}
After the Translator emits code, a second agent with a role-flipped system prompt (``you are a critical reviewer who must find at least one bug'') inspects the output; if it returns no defect, the Translator is re-prompted with a meta-instruction to re-check for common bug classes (ChatDev [P46] \S2.4).

\item \textbf{Symbol-Aware Navigator.}
A dedicated retrieval agent that exposes \texttt{SearchSymbols}, \texttt{GetDefinition}, \texttt{GetReferences}, and \texttt{SearchRepoGraph(term)}; before each fragment translation, the Translator queries the Navigator and receives a tight context window ($\le 4$\,KB) of code excerpts ranked by proximity and symbol importance (HyperAgent [P47] \S6.2 Table 6; RepoGraph [P62]).

\item \textbf{Coverage-Guided Plateau Detection.}
Combine search-based software testing with LLM-suggested tests injected when coverage plateaus are detected; halts unproductive repair loops when test gains stagnate (CodaMOSA [P49] ICSE 2023).

\item \textbf{Conversable State Checkpoints \& Interrupts.}
A translation pipeline can be paused mid-flight after each stage or iteration and resumed from persistent disk snapshots, enabling human review or external interrupt handling (AutoGen [P48]).

\item \textbf{Recruitment-Adaptive Planning.}
Replace fixed pipeline structures with a dynamic recruitment step before each iteration: pick per-iteration tool/agent configurations and re-evaluate after the validator pass (AgentVerse [P26] \S2).

\item \textbf{Source-to-Target Manifest Rewriter.}
For each translation target, take the source dependency list, apply a curated crate mapping table, and emit a minimal target manifest containing only required crates; unmapped dependencies are isolated for explicit review (Syzygy [P58] \S5.3; JavaC2Rust [P60] \S3).

\item \textbf{Iterative Retrieval Refinement.}
At translation iteration $i$, the Navigator re-indexes with prior translated fragments to produce an updated, tighter context window for fragment $i+1$ (RepoCoder [P61]).

\end{enumerate}

\paragraph{Implementation Mapping.}
Key primitives implemented in the MAgHARCM codebase include:
\begin{itemize}
\item Reverse Topological Scheduling \& Back-Edge Cycle Removal --- \texttt{internal/agents/planning.go}
\item Multi-Stage Build/Test Feedback Repair --- \texttt{internal/agents/validator.go}
\item Adversarial Test-Weakening Guard --- \texttt{internal/agents/validator.go}
\item Migration Strategy Selection --- \texttt{internal/agents/analyzer.go}
\item Chunked Translation Dispatch --- \texttt{internal/agents/chunked\_translator.go}
\item Symbol-Aware Navigator --- \texttt{internal/agents/navigator.go}
\item Coverage Plateau Detection --- \texttt{internal/agents/plateau.go}
\item Conversable State Checkpoint Persistence --- \texttt{internal/agents/checkpoint.go}
\item Source-to-Target Manifest Rewriter --- \texttt{internal/agents/manifest_rewriter.go}
\end{itemize}
